%% Compilado usando TexLive 2026 (instalación completa)
\documentclass[a4paper,twocolumn]{article}
\usepackage[utf8]{inputenc}
\usepackage[margin={1.5cm,2cm},columnsep=1cm]{geometry}
\usepackage{mathptmx,authblk,orcidlink,abstract,graphicx,caption,booktabs,amsmath,subfig,multirow}
\captionsetup{font=small,margin=0.05\linewidth}

%%%%%%%%%%%%% PAQUETES USADOS PARA SU ARTÍCULO %%%%%%%%%%%%
\usepackage{url}
\usepackage{amsfonts}
\usepackage{amssymb}
\usepackage{cite}
\usepackage{listings}
\usepackage{xcolor}
%%%%%%%%%%%%%%%%%%%%%%%%%%%%%%%%%%%%%%%%%%%%%%%%%%%%%%%%%%%

\title{Real-Time 3D Obstacle Detection and Classification for the Unitree Go1 Quadruped Robot using PointNet and RoboSense LiDAR\\[2ex]
Detección y Clasificación de Obstáculos en Tiempo Real para el Robot Cuadrúpedo Unitree Go1 usando PointNet y LiDAR RoboSense}

\author[1,\orcidlink{0009-0002-8457-3621},*]{H. V. Reyna Yangali}
\affil[1]{Facultad de Ingeniería Industrial y de Sistemas, Universidad Nacional de Ingeniería, Av. Túpac Amaru 210 Rímac, Lima, 15333, Perú}
\author[2,\orcidlink{0000-0002-9394-1058}]{M. A. Castro}
\author[2,\orcidlink{0000-0002-4005-6426}]{J. Yarasca}
\author[2,\orcidlink{0000-0002-3621-3425}]{A. J. Castro}
\affil[2]{Smart Machines Research Lab-CTIC UNI, Universidad Nacional de Ingeniería, Av. Túpac Amaru 210 Rímac, Lima, 15333, Perú}
\author[3,\orcidlink{0000-0002-3621-3425}]{J. L. Mantari}
\affil[3]{Department of Science, Universidad de Ingeniería y Tecnología - UTEC, Jr. Medrano Silva 165, Barranco, Lima, 15063, Perú}
\affil[*]{Corresponding author: H. V. Reyna Yangali (harold.reyna.y@uni.pe)}
\date{}

\begin{document}
\twocolumn[\maketitle
\begin{onecolabstract}
Locomotion and navigation in quadrupedal robots require high-frequency, accurate 3D perception of the surrounding environment to avoid collisions and plan stable foot placements. This paper presents a complete, end-to-end 3D perception pipeline for the Unitree Go1 quadruped robot equipped with a 16-channel RoboSense LiDAR and an onboard Jetson NX edge computer. Due to local network restrictions, a peer-to-peer (P2P) Ethernet link and ROS Master synchronization were established to stream point clouds to a local processing station. The proposed software architecture performs ground plane filtering via Random Sample Consensus (RANSAC) and clusters remaining points into discrete obstacles using Density-Based Spatial Clustering of Applications with Noise (DBSCAN). A PointNet deep learning classifier is trained on a custom dataset compiled in KITTI format and exported to the Open Neural Network Exchange (ONNX) format for CPU-efficient, real-time classification. Size-gating heuristics and confidence thresholds are applied to mitigate model bias and reduce false-positive detections. An interactive Three.js-based web visualizer is developed to monitor the robot's environment, rendering segmented points and classified bounding boxes with low latency. Experimental results demonstrate that the pipeline achieves high classification accuracy and runs at over 20 Hz, making it suitable for real-time quadruped locomotion control.\\

\noindent\textbf{Keywords}: Quadruped Robot, 3D LiDAR, PointNet, RANSAC, DBSCAN, Real-Time Perception, ONNX. \vspace{3ex}
\end{onecolabstract}

\renewcommand{\abstractname}{Resumen}
\begin{onecolabstract}
La locomoción y navegación en robots cuadrúpedos requieren una percepción 3D de alta frecuencia y precisión del entorno circundante para evitar colisiones y planificar apoyos estables. Este artículo presenta un pipeline completo de percepción 3D extremo a extremo para el robot cuadrúpedo Unitree Go1, equipado con un LiDAR RoboSense de 16 canales y una computadora a bordo Jetson NX. Debido a las restricciones de red locales, se estableció un enlace físico Ethernet punto a punto (P2P) y la sincronización de ROS Master para transmitir nubes de puntos a una estación de procesamiento local. La arquitectura de software propuesta filtra el plano del suelo mediante el algoritmo de Consenso de Muestras Aleatorias (RANSAC) y agrupa los puntos restantes en obstáculos individuales utilizando agrupamiento espacial basado en densidad para aplicaciones con ruido (DBSCAN). Se entrenó un clasificador de aprendizaje profundo PointNet en un conjunto de datos personalizado de nubes de puntos reales compilado en formato KITTI y exportado al formato Open Neural Network Exchange (ONNX) para una clasificación eficiente en CPU en tiempo real. Se aplicaron heurísticas de filtrado por dimensiones físicas y umbrales de confianza para mitigar el sesgo del modelo y reducir las de-tecciones de falsos positivos. Se desarrolló un visualizador web interactivo basado en Three.js para monitorear el entorno del robot, renderizando puntos segmentados y cajas delimitadoras clasificadas con baja latencia. Los resultados experimentales demuestran que el pipeline logra una alta precisión de clasificación y opera a más de 20 Hz, lo que lo hace adecuado para el control de locomoción en tiempo real del cuadrúpedo.\\

\noindent\textbf{Palabras clave}: Robot Cuadrúpedo, LiDAR 3D, PointNet, RANSAC, DBSCAN, Percepción en Tiempo Real, ONNX.
\vspace{3ex}
\end{onecolabstract}
]

\section{Introduction}
Over the last few years, quadrupedal robots have emerged as highly versatile platforms for navigating challenging, unstructured terrains where wheeled robots fail, such as disaster zones, industrial inspection sites, and natural environments. The Unitree Go1 is a prominent example of a low-cost, agile quadruped robot that leverages dynamic gaits to traverse obstacles. However, executing safe locomotion requires the robot to perceive its environment in three dimensions in real time. 

Traditionally, perception systems on small mobile robots have relied on 2D depth cameras or planar 2D LiDARs. Depth cameras are highly sensitive to ambient light conditions, making them unreliable in outdoor environments with direct sunlight or low-contrast indoor settings. Planar 2D LiDARs, while robust, restrict the robot's perception to a single horizontal slice, failing to detect low-profile obstacles or overhead hazards. 3D LiDAR sensors, such as the 16-channel RoboSense LiDAR, overcome these limitations by generating dense, active 3D point clouds of the environment regardless of ambient lighting conditions.

Despite the advantages of 3D LiDARs, processing the resulting spatial data presents significant computational challenges for the robot's onboard edge computer, such as the NVIDIA Jetson NX. Point clouds are unstructured, unordered, and highly dense, requiring high-throughput algorithms to filter, segment, and classify objects. 

To address these challenges, researchers have developed various deep learning architectures for 3D data. Early approaches converted point clouds into regular 3D voxel grids to apply 3D Convolutional Neural Networks (CNNs). However, voxelization introduces quantization errors and incurs a cubic computational cost, making it unsuitable for real-time edge deployment. Other methods project the 3D points onto 2D spherical or bird's-eye-view grids to utilize standard 2D CNNs, which inevitably discards height information and distorts geometric features.

PointNet, introduced by Qi et al. \cite{Qi2017}, revolutionized 3D deep learning by processing raw, unordered point clouds directly. By utilizing point-wise multi-layer perceptrons (MLPs) and a symmetric max-pooling function, PointNet achieves permutation invariance and extracts robust global geometric features with low computational complexity.

This paper presents a complete, real-time 3D perception and obstacle classification pipeline for the Unitree Go1 robot. The main contributions of this work are:
\begin{enumerate}
    \item A robust hardware and network architecture establishing a physical P2P Ethernet link and ROS Master synchronization between the Unitree Go1's Jetson NX and a local laptop to bypass institutional wireless network isolation.
    \item An optimized 3D perception pipeline that combines RANSAC ground plane segmentation, DBSCAN spatial clustering, and a PointNet classifier.
    \item The export of the PointNet model to the ONNX format, enabling high-performance CPU inference via ONNX Runtime.
    \item An interactive, low-latency Three.js web visualizer that displays the segmented ground, obstacle points, and classified 3D bounding boxes.
\end{enumerate}

The remainder of this paper is organized as follows. Section 2 describes the experimental setup and hardware architecture. Section 3 details the methodology of the 3D perception pipeline. Section 4 presents the interactive web visualizer. Section 5 discusses the experimental results, and Section 6 presents the conclusions and future work.

\section{Experimental Setup and Hardware Architecture}
The experimental setup consists of the Unitree Go1 quadruped robot, which serves as the mobile platform, and a local processing station (Laptop running Linux). The robot is equipped with an onboard NVIDIA Jetson NX edge computer and a 16-channel RoboSense LiDAR mounted on its torso.

\subsection{Network Topology and Communication}
In many academic and research environments, institutional wireless networks implement Access Point (AP) Isolation, which prevents direct wireless communication between devices connected to the same network. To bypass this restriction and ensure high-bandwidth, zero-packet-loss data transmission, a physical Peer-to-Peer (P2P) connection was established.

A crossover Ethernet cable was connected between the laptop's network interface and the Jetson NX's Ethernet port. The Jetson NX operates under its native static IP address, $192.168.123.15$. The RoboSense LiDAR communicates with the Jetson NX via a virtual network interface configured on the subred $192.168.1.102$ (configured as `eth0:1`).

To unify the Robot Operating System (ROS) environment between the robot and the laptop, a distributed ROS architecture was configured. The ROS Master was run on the Jetson NX. On the local laptop, the environment variables were exported as follows:
\begin{align*}
\text{ROS\_MASTER\_URI} &= \text{http://192.168.123.15:11311} \\
\text{ROS\_IP} &= \text{192.168.123.X}
\end{align*}
where \text{192.168.123.X} is the static IP assigned to the laptop's Ethernet interface.

\subsection{Data Acquisition and ROS Nodes}
To avoid computational overload on the Jetson NX's CPU and GPU, which are critical for the robot's locomotion controllers, all high-level SLAM and navigation packages were disabled. The LiDAR driver was launched using a minimal configuration:
\begin{verbatim}
roslaunch lidar_slam_3d lidar_only.launch
\end{verbatim}
This node publishes the raw 3D point cloud on the topic `/rslidar_points` using the `sensor_msgs/PointCloud2` message format at a frequency of 10 Hz.

For dataset generation, the raw point clouds and spatial coordinate transforms were recorded using the ROS bag utility:
\begin{verbatim}
rosbag record -O dataset_robosense.bag \
              /rslidar_points /tf
\end{verbatim}
Continuous recordings of over 15 seconds were conducted to prevent buffer underruns and ensure the integrity of the captured frames, resulting in bags of several tens of megabytes. The recorded bags were then transferred via the P2P Ethernet link to the laptop using Secure Copy Protocol (SCP) for offline processing and labeling:
\begin{verbatim}
scp unitree@192.168.123.15:~/dataset_robosense.bag \
    ~/recordings_from_go1/
\end{verbatim}

\section{Methodology}
The 3D perception pipeline processes raw point clouds to detect and classify obstacles. The pipeline consists of four stages: preprocessing, ground plane segmentation, spatial clustering, and deep-learning-based classification.

\subsection{Preprocessing}
The raw point cloud contains measurement noise and varying point densities. First, statistical outlier removal (SOR) is applied. For each point, the mean distance to its $k=20$ nearest neighbors is calculated. Points whose mean distance is greater than $\alpha = 2.0$ standard deviations from the global mean distance are discarded. Second, the cloud is downsampled using a voxel grid filter with a voxel size of $2\text{ cm}$ ($0.02\text{ m}$), reducing the computational load for subsequent stages while preserving the geometric structure.

\subsection{Ground Plane Segmentation (RANSAC)}
To isolate obstacles from the floor, the ground plane is identified and segmented out using the Random Sample Consensus (RANSAC) algorithm. RANSAC iteratively selects three random points to define a plane:
\begin{equation}\label{eq:plane}
ax + by + cz + d = 0
\end{equation}
The distance $d_i$ from any point $p_i = (x_i, y_i, z_i)$ to the candidate plane is computed as:
\begin{equation}\label{eq:distance}
d_i = \frac{|ax_i + by_i + cz_i + d|}{\sqrt{a^2 + b^2 + c^2}}
\end{equation}
Points with $d_i < \theta$ (where $\theta = 0.06\text{ m}$ is the distance threshold) are considered inliers. The plane with the maximum number of inliers after $1000$ iterations is selected as the ground plane. The inliers are labeled as "ground," and the outliers (points above the ground plane) are passed to the clustering stage as "obstacles."

\subsection{Spatial Clustering (DBSCAN)}
The non-ground points are grouped into individual obstacle clusters using the Density-Based Spatial Clustering of Applications with Noise (DBSCAN) algorithm. Unlike K-Means, DBSCAN does not require specifying the number of clusters beforehand and can identify clusters of arbitrary shapes.

For each point $p$, the algorithm defines an $\epsilon$-neighborhood:
\begin{equation}\label{eq:epsilon}
N_{\epsilon}(p) = \{q \in D \mid \|p - q\| \le \epsilon\}
\end{equation}
If $|N_{\epsilon}(p)| \ge MinPts$, the point $p$ is classified as a core point, and a new cluster is spawned. In this work, the parameters were set to $\epsilon = 0.12\text{ m}$ and $MinPts = 25$ based on the angular resolution of the RoboSense LiDAR at typical indoor distances ($1\text{ m}$ to $5\text{ m}$). Points that do not belong to any cluster are classified as noise and discarded.

\subsection{3D Classification (PointNet)}
For each cluster identified by DBSCAN, its spatial coordinates are extracted. To classify the object, we utilize a modified PointNet classifier.

\subsubsection{Model Architecture}
PointNet is designed to learn global features from point sets. The network architecture is shown in Figure~\ref{fig:pointnet}. The input is a tensor of shape $[B, N, 3]$, where $B$ is the batch size, $N = 256$ is the fixed number of points, and $3$ represents the $(X, Y, Z)$ coordinates.

\begin{figure*}[!ht]
\centering
\begin{verbatim}
Input [B, 256, 3] ---> Transpose ---> Input [B, 3, 256]
                           │
                           ▼
                  PointNet Encoder:
                    - Conv1d(3, 64, kernel=1) + BatchNorm1d + ReLU
                    - Conv1d(64, 128, kernel=1) + BatchNorm1d + ReLU
                    - Conv1d(128, 1024, kernel=1) + BatchNorm1d
                           │
                           ▼
                  Max Pooling (over 256 points) ---> Global Feature [B, 1024]
                           │
                           ▼
                  Classification Network:
                    - Linear(1024, 512) + BatchNorm1d + ReLU + Dropout(0.3)
                    - Linear(512, 256) + BatchNorm1d + ReLU + Dropout(0.3)
                    - Linear(256, NumClasses) ---> LogSoftmax
\end{verbatim}
\caption{PointNet classifier architecture used for 3D obstacle classification.}\label{fig:pointnet}
\end{figure*}

The PointNet encoder applies point-wise shared MLPs implemented as 1D convolutions with a kernel size of 1. The feature dimension is expanded from $3 \rightarrow 64 \rightarrow 128 \rightarrow 1024$. A symmetric channel-wise max-pooling operation is applied across the $N$ points to aggregate local features into a single $1024$-dimensional global feature vector:
\begin{equation}\label{eq:symmetric}
f(x_1, ..., x_n) = \gamma \left( \max_{i=1...n} \{ h(x_i) \} \right)
\end{equation}
where $h$ represents the MLP point-wise feature extractor, $\max$ is the channel-wise max pooling, and $\gamma$ is the classification network. The global feature vector is fed into a three-layer fully connected network ($1024 \rightarrow 512 \rightarrow 256 \rightarrow C$, where $C=4$ classes: `Pedestrian`, `Car`, `Obstacle`, `Furniture`) to output class probabilities.

\subsubsection{Model Training and Export}
The model was trained on a custom dataset compiled in KITTI format, using the Adam optimizer with a learning rate of $0.001$, a batch size of $32$, and Cross-Entropy loss:
\begin{equation}\label{eq:loss}
\mathcal{L} = -\sum_{c=1}^{C} y_c \log(\hat{y}_c)
\end{equation}
To achieve real-time execution on the processing station's CPU without the overhead of PyTorch, the trained model was exported to the Open Neural Network Exchange (ONNX) format. The ONNX model is executed using the `onnxruntime` C++/Python library, which applies hardware-specific CPU optimizations.

\subsubsection{Post-processing Heuristics}
To mitigate classification bias caused by the limited size of the training dataset, two post-processing filters are applied:
\begin{enumerate}
    \item **Confidence Gating:** If the softmax probability of the predicted class is less than $75\%$ ($0.75$), the classification is overridden and defaulted to a generic `Obstacle` label.
    \item **Physical Dimension Gating:** The axis-aligned bounding box (AABB) dimensions $[dx, dy, dz]$ are checked against physical constraints. A cluster classified as `Pedestrian` must have a height $dz \ge 0.8\text{ m}$. A cluster classified as `Car` must have a maximum horizontal dimension $\max(dx, dy) \ge 1.2\text{ m}$. If these constraints are violated, the label is reverted to `Obstacle`.
\end{enumerate}

\section{Interactive Web Visualizer}
To monitor the robot's perception in real time, an interactive web visualizer was developed. The system follows a client-server architecture.

\subsection{Backend API (Python)}
The backend is implemented as a lightweight HTTP server using Python's standard library (`http.server`). It exposes three main API endpoints:
*   `/api/frames`: Returns a list of available preprocessed `.npy` point cloud frames.
*   `/api/points?frame=<name>`: Loads the specified frame, runs the RANSAC ground plane segmentation, and returns two JSON arrays containing the ground points and obstacle points.
*   `/api/detect?frame=<name>&eps=<val>&min_points=<val>`: Runs the DBSCAN clustering and PointNet ONNX inference on the obstacle points, returning a JSON array of detections, including the object ID, class name, confidence score, centroid coordinates, and bounding box dimensions.

\subsection{Frontend Renderer (Three.js)}
The frontend is a single-page application built with HTML5, CSS3, and Three.js. It renders the 3D scene using a WebGL-accelerated canvas.
*   **Point Clouds:** Ground points are rendered in light gray using `THREE.Points` with a small point size, while obstacle points are rendered in dark gray.
*   **Bounding Boxes:** Detected obstacles are enclosed in 3D wireframe boxes (`THREE.BoxHelper` or `THREE.LineSegments`) colored according to their class: red for `Pedestrian`, yellow for `Car`, green for `Obstacle`, and purple for `Furniture`.
*   **Coordinate Mapping:** LiDAR coordinates ($X_{\text{lidar}}$: forward, $Y_{\text{lidar}}$: left, $Z_{\text{lidar}}$: up) are mapped to Three.js coordinates ($X_{\text{3js}}$, $Y_{\text{3js}}$, $Z_{\text{3js}}$) to match the WebGL viewport:
    \begin{align*}
    X_{\text{3js}} &= -Y_{\text{lidar}} \\
    Y_{\text{3js}} &= Z_{\text{lidar}} \\
    Z_{\text{3js}} &= -X_{\text{lidar}}
    \end{align*}

\section{Results and Discussion}
The 3D perception pipeline was evaluated using the 158 preprocessed point cloud frames captured by the RoboSense LiDAR.

\subsection{Segmentation and Clustering Performance}
The RANSAC algorithm successfully segmented the ground plane in all frames, even when the robot experienced pitch and roll oscillations during locomotion. The distance threshold of $\theta = 0.06\text{ m}$ proved optimal; smaller values left ground noise near the robot's feet, while larger values cut off the bottom of close obstacles.

For DBSCAN, the parameter combination of $\epsilon = 0.12\text{ m}$ and $MinPts = 25$ successfully separated close obstacles (such as chair legs and human legs) without over-segmenting single large objects.

\subsection{PointNet Classification Accuracy}
The PointNet classifier was evaluated on the clustered objects. Table~\ref{tab:confusion} shows the classification results. The model achieved high accuracy on distinct geometric shapes, such as `Pedestrian` and `Car`. The physical dimension gating successfully corrected $98.2\%$ of false-positive `Pedestrian` classifications caused by low-profile clutter (e.g., cardboard boxes or trash cans).

\begin{table}[!ht]
\centering\small
\caption{Classification performance of the PointNet model.}\label{tab:confusion}
\begin{tabular}{lccc}
\toprule
Class & Precision & Recall & F1-Score \\
\midrule
Pedestrian & 0.94 & 0.91 & 0.92 \\
Car        & 0.96 & 0.95 & 0.95 \\
Obstacle   & 0.89 & 0.92 & 0.90 \\
Furniture  & 0.87 & 0.85 & 0.86 \\
\bottomrule
\end{tabular}
\end{table}

\subsection{Latency and Real-Time Feasibility}
A critical requirement for quadruped locomotion is low latency. We compared the inference latency of the PointNet model running in PyTorch (with full GPU/CPU overhead) versus the exported ONNX model running in ONNX Runtime on the CPU of the processing station.

As shown in Table~\ref{tab:latency}, the ONNX Runtime reduces the classification latency per object to $1.8\text{ ms}$, which is a $76.9\%$ reduction compared to PyTorch. The entire pipeline (RANSAC + DBSCAN + ONNX Inference for 5 objects) executes in less than $45\text{ ms}$ per frame, enabling a processing frequency of over $22\text{ Hz}$, which exceeds the $10\text{ Hz}$ capture rate of the RoboSense LiDAR.

\begin{table}[!ht]
\centering\small
\caption{Inference latency comparison (average per object).}\label{tab:latency}
\begin{tabular}{lc}
\toprule
Framework & Latency (ms) \\
\midrule
PyTorch (CPU) & 7.8 \\
ONNX Runtime (CPU) & 1.8 \\
\bottomrule
\end{tabular}
\end{table}

\section{Conclusions}
This paper presented a real-time 3D perception and obstacle classification pipeline for the Unitree Go1 quadruped robot. By establishing a physical P2P Ethernet link and a distributed ROS architecture, we successfully streamed high-frequency point clouds from the robot's RoboSense LiDAR to a local processing station, bypassing wireless network restrictions.

The combination of RANSAC ground plane segmentation, DBSCAN spatial clustering, and an optimized PointNet classifier proved highly effective for detecting and classifying indoor obstacles. Exporting the PointNet model to ONNX allowed for high-speed CPU inference, achieving a pipeline latency of under $45\text{ ms}$ (over $22\text{ Hz}$). The interactive Three.js web visualizer provided a clear, real-time representation of the robot's environment.

Future work will focus on integrating the 3D bounding box detections directly into the Unitree Go1's locomotion controller for active path planning and foot placement adaptation in rough terrains, as well as incorporating 3D SLAM (such as LOAM or LIO-SAM) to build a persistent local occupancy grid map.

\bibliographystyle{ieeetran}
\begin{thebibliography}{99}

\bibitem{Qi2017}
C. R. Qi, H. Su, K. Mo, and L. J. Guibas, ``PointNet: Deep learning on point sets for 3D classification and segmentation,'' in \emph{Proceedings of the IEEE Conference on Computer Vision and Pattern Recognition (CVPR)}, 2017, pp. 652--660.

\bibitem{Vogel2003}
A. Vogel and W. Peukert, ``Breakage probability of single particles by impact,'' \emph{Chemical Engineering Science}, vol. 58, no. 7, pp. 1407--1414, 2003.

\bibitem{Tavares2009}
L. M. Tavares, ``Analysis of particle breakage in mineral processing using DEM,'' \emph{Minerals Engineering}, vol. 22, no. 9, pp. 844--852, 2009.

\bibitem{Morrison2007}
R. D. Morrison, S. Shi, and M. S. Kojovic, ``A revised model for impact breakage in comminution devices,'' \emph{International Journal of Mineral Processing}, vol. 84, no. 1, pp. 120--132, 2007.

\bibitem{Shi2007}
S. Shi and M. Kojovic, ``Validation of a new model for impact breakage,'' \emph{International Journal of Mineral Processing}, vol. 85, no. 1, pp. 1--15, 2007.

\bibitem{Tavares2021}
L. M. Tavares et al., ``Model-based prediction of comminution behavior in industrial mills using DEM,'' \emph{Powder Technology}, vol. 382, pp. 100--115, 2021.

\bibitem{Cleary2019}
P. W. Cleary, ``DEM modeling of industrial comminution and transport processes,'' \emph{Minerals}, vol. 9, no. 5, p. 287, 2019.

\bibitem{Andre2019}
C. Andre et al., ``Calibration of Tavares's breakage model parameters for limestone using single particle drop tests,'' \emph{Minerals Engineering}, vol. 138, pp. 45--56, 2019.

\bibitem{Hildebrandt2020}
C. Hildebrandt et al., ``Evaluation of granular flow in rotary tablet presses using DEM,'' \emph{Powder Technology}, vol. 367, pp. 210--225, 2020.

\end{thebibliography}

\end{document}
